% =====================================================================
% Section 3 -- Analysis
% Self-contained draft. Notation is defined from scratch; no dependency
% on any other note file. Citation keys reused from the existing
% Related Work section (niu2023sar, doloriel2026daf, xiong2026mgp,
% wang2022cotta, noori2025mlmp) -- add these entries to the .bib if not
% already present.
% =====================================================================

\section{Analysis} \label{analysis}

We analyze why continual adaptation collapses under long horizons, and use this
analysis to justify the two label-free signals introduced in
Section~\ref{intro}. Throughout, $\theta$ denotes the adapted parameters
(the visual-encoder LayerNorm affine parameters, following prior OVSS
adaptation work), and $p_\theta(\cdot \mid x) \in \Delta^{C-1}$ the predicted
class distribution for input $x$ over the $C$ classes present in the current
vocabulary, obtained by a softmax over visual--text compatibility scores.
Most CTTA objectives, including the OVSS baselines we compare against, adapt
$\theta$ by minimizing the expected prediction entropy
$L(\theta) = \mathbb{E}_x\big[H(p_\theta(\cdot\mid x))\big]$, where
$H(q) = -\sum_i q_i \log q_i$. Given a window of $N$ consecutive unlabeled
predictions during adaptation, we write their average as the
\emph{class-marginal distribution}
$\bar y = \frac1N\sum_{n=1}^N p_\theta(\cdot \mid x_n)$, with class-marginal
entropy $H(\bar y)$.

\subsection{Why does collapse happen?} \label{sec:why-collapse}

Entropy minimization rewards confident predictions but places no explicit
constraint on \emph{which} class that confidence falls on. Treating the
sample index as a random variable, the class-marginal entropy and the
per-sample conditional entropy decompose the mutual information between
inputs and predictions:
\begin{equation}
I(x; c) \;=\; H(\bar y) \;-\; \mathbb{E}_x\big[H(p_\theta(\cdot\mid x))\big].
\label{eq:mi-decomp}
\end{equation}
The adaptation objective $L(\theta)$ minimizes only the second term. Nothing
in $L(\theta)$ penalizes $H(\bar y) \to 0$: a model that maps \emph{every}
input to the same class $k$ attains $L(\theta) = 0$, the global minimum of the
objective, together with $H(\bar y) = 0$. In other words, the single-class
collapse pattern we observe empirically is not a failure of optimization --
it is the objective's own optimum.

A natural response, and one taken by prior OVSS work, is to add a
marginal-diversity term that explicitly rewards $H(\bar y)$, e.g.\ the
regularized objective used by DAF~\citep{doloriel2026daf},
$L_\lambda(\theta) = \mathbb{E}_x[H(p_\theta(\cdot\mid x))] - \lambda H(\bar
y_\theta)$. We show that this fix is structurally limited:

\begin{quote}
\textbf{Proposition 1.} \textit{Let $\{\theta^{(s)}\}$ be a sequence of
parameters along which predictions saturate toward a single class, i.e.\
$p_{\theta^{(s)}}(\cdot\mid x) \to e_k$ for (almost) every $x$ and a common
$k$. Then $\nabla_\theta L_\lambda(\theta^{(s)}) \to 0$ for every finite
$\lambda \in \mathbb{R}$.}
\end{quote}

\emph{Proof sketch.} Differentiating $H(p)$ with respect to the logits $z$
gives the closed form $\partial H(p)/\partial z_j = -p_j(\log p_j + H(p))$,
and an analogous expression for $\partial H(\bar y)/\partial z_j$ scaled by
$1/N$ (full derivation in Appendix~\ref{app:proofs}). Both expressions carry
an explicit factor of $p_j$ (or the corresponding per-sample probability),
which vanishes as the prediction saturates. The conditional-entropy term and
the diversity term therefore vanish \emph{at the same rate}, so their
weighted sum vanishes regardless of $\lambda$. $\blacksquare$

The mechanism is softmax saturation: once a prediction is sufficiently
confident, the model becomes locally insensitive to further changes in its
own output, and this insensitivity multiplies through every term built from
$H(p)$ or $H(\bar y)$. Consequently, a diversity term can slow the approach
to the degenerate solution -- and we confirm this empirically, where adding
such a term delays but does not prevent eventual collapse under our extended
horizon -- but it cannot alter the destination, because the degenerate
solution remains a critical point of the regularized objective for any
regularization strength. This gives a principled reason, rather than only an
empirical one, for why our method intervenes by \emph{restoring} parameters
rather than by adding another gradient-based term to the loss.

This collapse mechanism is a property of entropy minimization itself and does
not depend on the openness or size of the class vocabulary: it threatens any
method that adapts by minimizing prediction entropy, open-vocabulary or
closed-set alike, and is consistent with instability reported for
closed-set entropy-based TTA~\citep{niu2023sar}. It is, however, a different
mechanism from the one recently identified for long-horizon closed-set
classification by MGP~\citep{xiong2026mgp}, who attribute collapse to
spurious gradient components leaking into a protected parameter subspace and
correct it by projecting the entropy gradient's direction within that
subspace. Because this correction acts purely on the direction of the
gradient and not on the marginal distribution of predictions, it does not
address the class-marginal degeneracy described above -- consistent with our
observation that MGP itself becomes unstable under our long-horizon OVSS
setting (Section~\ref{sec:results}).

\subsection{Signal: gradient norm} \label{sec:signal-gradnorm}

The proof of Proposition~1 gives gradient norm a concrete interpretation.
Because $\partial H(p)/\partial z_j = -p_j(\log p_j + H(p))$ vanishes whenever
a prediction is confident, the norm of the adaptation gradient is small both
when the model has settled into a well-calibrated, confidently correct state
\emph{and} when it has collapsed onto a degenerate, confidently wrong one --
both are low-entropy, low-gradient configurations. Gradient norm at
convergence is therefore ambiguous between these two outcomes and cannot,
by itself, certify that a low-gradient state is a safe one to consider
reliable.

Away from either endpoint, however, gradient norm behaves informatively. As
the model transitions out of a stable configuration -- its predictions no
longer matching the confident, calibrated state that produced low gradients,
but not yet saturated onto a degenerate class either -- the entropy loss
surface is locally steeper, and the resulting gradient magnitude increases.
This transitional rise, consistent with gradient behavior reported during
the onset of instability in entropy-based TTA~\citep{niu2023sar}, is what
makes gradient norm useful as an \emph{early} indicator: it grows while the
model is drifting away from a reliable state, before predictions have fully
saturated onto a single class and the gradient vanishes again by
Proposition~1. Figure~\ref{fig:gradnorm} shows this pattern against mIoU
across our benchmarks: gradient norm rises as accuracy begins to degrade, and
only collapses toward zero once degradation is already severe -- i.e., it
reliably signals instability while there is still adaptation quality worth
protecting, but its magnitude alone cannot distinguish an already-collapsed
state from a genuinely healthy one. This motivates pairing it with a signal
that can make that distinction: class-marginal entropy.

\subsection{Signal: class-marginal entropy} \label{sec:signal-hmargin}

Class-marginal entropy directly targets the term that the adaptation
objective ignores (Equation~\ref{eq:mi-decomp}). Writing $u$ for the uniform
distribution over the $C$ classes,
\begin{equation}
H(\bar y) \;=\; \log C \;-\; D_{\mathrm{KL}}(\bar y \,\|\, u),
\label{eq:hmargin-kl}
\end{equation}
an exact identity (expand the KL definition and use $\sum_i \bar y_i = 1$).
Equation~\ref{eq:hmargin-kl} gives $H(\bar y)$ a direct information-geometric
reading: monitoring class-marginal entropy is equivalent, up to the constant
$\log C$, to monitoring how far the aggregated predictions have drifted from
uniform use of the vocabulary. Equivalently, writing the \emph{effective
number of classes in use} as $C_{\mathrm{eff}} = \exp(H(\bar y))$, a drop in
$H(\bar y)$ corresponds to a multiplicative shrinkage of $C_{\mathrm{eff}}$
toward $1$ -- a direct formalization of "probability mass becoming
concentrated on a small subset of categories."

A natural first bound, $H(\bar y) \ge -\log(\max_i \bar y_i)$, is tight but
uninformative at realistic operating points: on our benchmarks, marginal
entropy near the onset of collapse is still far from its degenerate limit,
and this bound only certifies that \emph{some} class holds a small fraction
of the mass -- true of many healthy distributions as well. A quantitatively
useful guarantee instead follows from bounding the KL divergence in
Equation~\ref{eq:hmargin-kl} by a $\chi^2$-divergence and relating it to total
variation distance (Appendix~\ref{app:proofs}):
\begin{equation}
\big\| \bar y - u \big\|_{\mathrm{TV}} \;\ge\; \frac12\sqrt{\frac{\log C -
H(\bar y)}{C}}.
\label{eq:tv-bound}
\end{equation}
Equation~\ref{eq:tv-bound} says that once $H(\bar y)$ has dropped by a given
amount from its running maximum, the marginal prediction distribution is
\emph{guaranteed} to have moved at least a corresponding, quantifiable
fraction of its probability mass away from a healthy, diverse baseline. This
is what licenses treating a relative drop in $H(\bar y)$, rather than its
absolute value, as a collapse trigger: the guarantee holds regardless of
$C$, so the same relative threshold is meaningful across datasets and
vocabularies of very different size (see Section~\ref{sec:results} for the
transfer of a single threshold across all three benchmarks).

Class-marginal entropy has a blind spot that mirrors gradient norm's: because
it is a \emph{marginal}, aggregate statistic, it stays high whenever
predictions are heterogeneous across the aggregation window, even if
individual predictions are already degrading in quality -- so long as the
errors are spread across, rather than concentrated onto, a common class. On
benchmarks where degradation manifests this way, $H(\bar y)$ can remain near
its healthy value while accuracy declines, giving no warning on its own.
This is precisely the failure mode that gradient norm, tracking the
optimization dynamics rather than the prediction distribution, remains
sensitive to. The two signals are therefore complementary rather than
redundant: gradient norm can catch degradation that does not concentrate
onto a single class, and class-marginal entropy can distinguish a
low-gradient collapsed state from a low-gradient healthy one -- the
ambiguity left open at the end of Section~\ref{sec:signal-gradnorm}. Section
\ref{sec:method} builds our restoration mechanism, CAGA, jointly on both.
